\subsection{Cross-domain extension and stronger graph controls}
We next screened additional public sources with bounded local storage, retaining
unsuccessful and unsuitable candidates. Azure V2 \citep{azure2019} supplies a
second compute provider. A first 49,231-row schema/prediction pilot was followed
by a sample retaining \emph{all released VMs} in subscriptions selected by a
fixed SHA-256 rule (first 32 bits modulo 20 equal zero). The sample is complete
within those released subscriptions, not the provider's full production
population. The target is VM deletion at 5/30/120/360 minutes; observed deployment
and subscription memberships supply two relations. Lifetime CPU summaries,
usage-derived VM categories, and deletion times are excluded from features.
Time-zero creations give an observed age lower bound; timestamp coarsening and
the publisher's workload sampling limit interpretation. The interval adapter
emits creation events only, so Azure temporal tokens summarize prior creations,
not a full stream of VM deletions or CPU measurements.

The UCI incident log \citep{amaral2018} supplies handler and team relations.
We predict the first \emph{recorded} Resolved/Closed update after the first
observed nonterminal update at 1/6/24/72 hours. This is an event-log milestone,
not an estimate of an unobserved physical repair time. Resolution timestamps,
closure codes, eventual resolvers, and SLA outcomes repeated in earlier rows
are excluded. BPI 2014 \citep{vandongen2014activity} supplies timestamped incident
activities and team membership; its target is the first Closed activity after
an observed Open. We do not reconstruct live arrival populations from either
historical case extract. Event timestamps proxy information availability.
Naive incident timestamps are interpreted on a UTC clock for reproducible
partitioning; the archives do not verify operational timezone or ingestion
latency. This convention does not recover daylight-saving transitions.

All sources have distinct individual-fit, calibration, dependence-fit, and
development-validation windows. Entire longest-horizon labels must mature
strictly before each window ends. Azure uses day boundaries 0/10/13/16/20;
UCI uses March 1/22, April 1/15, and May 1, 2016; BPI 2014 uses October 1,
November 1/15, December 1, 2013, and January 1, 2014. Prefix state carries
across boundaries. New-period confirmation fixes 28 origins per source:
Azure days 22--28.75, UCI May 1--28, and BPI 2014 January 1--28.
An executable asset manifest freezes all fitted models and comparisons before
these targets are prepared or scored. Earlier schema audits inspected whole-file
state totals and time ranges; this is an internal prediction-holdout freeze,
not external preregistration or a claim that the raw files were never opened.

\paragraph{Matched individual models.}
Controls include personal-feature HGB, HGB with current relation summaries,
and HGB with the same observed event-history tokens available to temporal
attention. Neural controls use width 48: a node MLP, a global set model,
two rounds of typed entity--group mean message passing, and temporal attention
over the last eight observed relation events followed by message passing.
These are independent survival-target adaptations, not reproductions of TGN,
TGAT, or DyGFormer benchmark systems \citep{rossi2020,xu2020,yu2023}.
Each network uses a five-category first-exit head. UCI trains for 30 epochs;
Azure and BPI 2014 for 20, with AdamW learning rate .001 and weight decay .001.
HGB uses 100 boosting iterations, 15 leaves and L2 penalty 10. Checkpoints and temperature
are chosen on the earlier calibration window; three-seed mean CDFs determine
the selected architecture. Relation-specific and size-preserving shuffled
graph controls keep the input feature set and training budget fixed.
A separate bounded Borg pilot shares the existing 176 engineered features
across models, testing additional message passing beyond those summaries;
its sampled-node scores are not population estimates.

\paragraph{Joint controls.}
Gaussian factors include independent, global, each relation alone,
global-plus-each-relation, combined typed, and shuffled-incidence controls.
Small development-only variance grids and coordinate search fit count scores;
the typed model may reduce to the global control. Coordinate fitting starts
at the global fit and performs two sweeps over weights; it is not an exhaustive
joint-grid optimum. Separately fitted relation-only arms are retained even
when they beat that combined fit. Student-$t$ copulas add a
shared inverse-chi-square scale and compare global and typed variants with
degrees of freedom chosen from 4/8/16. Each draw still uses one entity uniform
across horizons. Concentration-only and bias-adjusted direct DM, exact
beta-binomial count scores, and historical residual forecasts use the same
dependence-fitting period. None is tuned on confirmation.

A small additional comparator uses the unchanged published TACTiS attentional
copula module \citep{drouin2022}, with a custom current-cohort encoder,
randomized discrete probability-integral-transform targets, and rank-adjusted
generated uniforms. All comparators in that pilot use the same 64-node subsets.
It is explicitly a limited adaptation: neither the full time-series model nor
an exact discrete likelihood reproduction, and not evidence of superiority
over TACTiS or TACTiS-2 in their original tasks.

Exploratory fitting uses empirical-ensemble CRPS; confirmation uses the fair
distinct-pair estimator with 512 draws and three sampling seeds. The peak
reference scale in confirmation is $\sqrt{1+\sum_iF_i(1-F_i)}$, as in the
original study. Development pilot peak scores used a max-with-one stabilization
and are not numerically pooled with confirmation. Seeds are averaged within
origin before 2,000 paired circular block replicates with 1/3/7-day blocks.
The same frozen probabilities and reference risk quantities are used for
every joint comparator.

\paragraph{Screened alternatives and limits.}
A Google 2019 \citep{google2019} event-shard audit found approximately 83\% of
sampled entity histories contained one event; that shard cannot support a
complete-trajectory forecast. BPI 2014 change snapshots lack version histories
for planned times and affected configuration items; incident activities offer
a better-supported relation replay. Citi Bike \citep{citibike} and airline
\citep{bts2024} pilots test transportation at small scale. Published trips are
retrospectively filtered, and airline schedule-version availability is an
explicit assumption. Their exploratory scores are not presented as live
deployment results. Source selection followed development feasibility and
performance, so reporting all screened candidates is part of the evidence.
